\documentclass{article}
\usepackage{geometry}
\usepackage{hyperref}
\usepackage{amsfonts}
\usepackage{amsmath}
\usepackage{graphicx}
\usepackage{/home/user1/code/tex/sty}
\setlength{\parindent}{0pt}
\geometry{top=1in,
  bottom=1in,
  left=1in,
  right=1in,
}
\newcommand{\pic}[2]{\begin{center}\includegraphics[height=#1]{#2}\end{center}}
\newcommand{\abs}[1]{|#1|}
\usepackage{graphicx}
\n{lorentz}{(1-v^2/c^{2})^{-1/2}}
\n{sxx}{\sigma_{xx}}
\n{syy}{\sigma_{yy}}
\n{sxy}{\sigma_{xy}}
\n{syx}{\sigma_{yx}}

\begin{document}
\begin{center}
  {\Large{Homework 3}}\\~\\
  \begin{tabular}{l l}
    Student & Agustin Romero\\
    Instructor & Sri Raghu\\
    Assistants & Josephine Yu, Pavel Nosov\\
    Course & Classical Electrodynamics\\
    Institution & Stanford University\\
    Due Date & 17 February 2023\\
  \end{tabular}
\end{center}
Collaborators:
\beg{enumerate}{
  \item Charles Yang
}
References:
\beg{enumerate}{
\item Griffiths p553
\item Landau Lifshitz p68
}

\section{Question 1.a.}
The magnetic field produced by both fields is 0.
\e{\vec{B}=0}
The electric field produced by the top plate is
\e{\abs{\vec{E}}=\fr{1}{2\epsilon_0}\sigma}
pointing away from the top plate. The bottom plate produces the field
\e{\abs{\vec{E}}=\fr{1}{2\epsilon_0}\sigma}
pointing towards the bottom plate. The two fields do not cancel between eachother, and the total field is
\e{\boxed{\vec{E}=0\hat{x}+\fr{1}{\epsilon_0}\sigma\hat{y}+0\hat{z}}}

\section{Question 1.b.i.}
\pic{5cm}{1.jpg}
Gauss' equation,
\e{\int \vec{E}d\vec{a}=\fr{1}{\epz}Q_{\rom{inc}}}
is valid for charges with constant velocity. In the rocket frame, Gauss' equation implies $\vec{E}$ at all coordinates between the two plates points along the \fbox{$-\hat{y}$} direction, with no $\hat{z}$ nor $\hat{x}$ components. Next, consider the $\vec{B}$ field. Ampere's equation is
\e{\nabla\times \vec{B}=\mu_0 \vec{J}}
The rocket observes that $\vec{J}$ is nonzero in the $-\hat{x}$ direction. By Ampere's equation, $\vec{B}$ points in the \fbox{$+\hat{z}$ direction}.

\section{Question 1.b.ii.}
Let $O'$ be the frame moving parallel to the plates, and $O$ be the plate frame.
\e{\boxed{\vec{E'}=0\hat{x}+\lorentz \fr{\sigma}{\epz}\hat{y}+0\hat{z}}}
\e{\boxed{\vec{B'}=0\hat{x}+0\hat{y}+\lorentz \muz \sigma v\hat{z}}}

\section{Question 1.c.i.}
Let $O'$ be the frame moving perpendicular to the plates, and $O$ be the plate frame. The rocket moves in the $\hat{y}$ direction to one plate. In the new frame, $\vec{E'}$ at all coordinates between the two plates points along the \fbox{$-\hat{y}$} direction, with no $\hat{z}$ nor $\hat{x}$ components. Next, consider the $\vec{B'}$ field. The rocket observes one plate moving towards it with velocity $-\vec{v}\hat{y}$, and the other plate moving away with velocity $+\vec{v}\hat{y}$. The rocket observes that neither plate generates a surface current $\vec{j}$, therefore Ampere's equation is
\e{\nabla \times \vec{B'} = 0}
and \fbox{$\vec{B'}=0$}.

\section{Question 1.c.ii.}
In the boosted frame, the $\vec{E}$ field is nonzero in the direction towards the plates, $\hat{y}$, but the $\vec{B}$ fields is zero.
\e{\boxed{\vec{E'}=0\hat{x}+\lorentz \fr{\sigma}{\epz}\hat{y}+0\hat{z}}}
\e{\boxed{\vec{B'}=0\hat{x}+0\hat{y}+0\hat{z}}}

\section{Question 1.d.}The boost equations a
\e{E'_{||}=E_{||}}
\e{B'_{||}=B_{||}}
\e{E'_\perp=\lorentz (E_\perp+c^{-1}\vec{v}\times\vec{B})}
\e{B'_\perp=\lorentz (B_\perp-c^{-1}\vec{v}\times\vec{E})}
Let $\vec{v}$ be parallel to the $\hat{x}$ axis. The component parallel to the boost, $\hat{x}$, is not transformed by the boost. For $E'_x$ to be 0, $E_x$ must be 0. The $\hat{y}$ and $\hat{z}$ components of $\vec{E}$ and $\vec{B}$ mix, and must be exactly negative of the resulting boosted fields $\vec{E'}$ and $\vec{B'}$. $\vec{v}\times \vec{B}$ is
\e{\vec{v}\times \vec{B}&=(v_yB_z-B_yv_z)\hat{x}-(v_xB_z-B_xv_z)\hat{y}+(v_xB_y-B_xv_y)\hat{z}\\
&=-v_xB_z\hat{y}+v_xB_y\hat{z}}
The condition is
\e{\label{noqwid}\boxed{\vec{E}=0\hat{x}+\fr{v_x}{c}B_z\hat{y}-\fr{v_x}{c}B_y\hat{z}}}

\section{Question 1.e.}The logic is identical to 1.d.. $\vec{v}\times \vec{E}$
\e{\vec{v}\times \vec{E}&=(v_yE_z-E_yv_z)\hat{x}-(v_xE_z-E_xv_z)\hat{y}+(v_xE_y-E_xv_y)\hat{z}\\
&=-v_xE_z\hat{y}+v_xE_y\hat{z}}
The condition is
\e{\label{aoduna}\boxed{\vec{B}=0\hat{x}-\fr{v_x}{c}E_z\hat{y}+\fr{v_x}{c}E_y\hat{z}}}

\section{Question 1.f.}
\pic{5cm}{2}
The Hall experiment can be summarized as:
\beg{enumerate}{
\item Setup two magnets, a line of current, and a dielectric as shown in figure.
\item Send current through the dielectric.
\item Observe that the current is temporarily perturbed, until a potential difference V appears in the $\hat{z}$ direction, and the current returns to its unperturbed path.
}
In this setup, $\vec{E}$ points in the $\hat{x}$ direction, $\vec{B}$ in the $\hat{y}$ direction. The Lorentz force implies that the current in the dielectric is forced in the $+\hat{z}$ direction
\e{\vec{E}\cdot\vec{B}=0}
and so the results of 1.d. and 1.e should apply. Observe that 1.d. predicts that \eref{noqwid} must be satisfied for there to exist a frame where $\vec{E}$ is zero, and 1.e. predicts that \eref{aoduna} must be satisfied for there to exist a frame where $\vec{B}$ is zero. Observe that the Hall experiment can satisfy either of these equations, where $v_x$ is the velocity of the electrons.

\section{Question 2.a.}
\e{\vec{B}=\vec{B'}-\fr{1}{c}\vec{E'}\times \vec{v}}
\e{\vec{B}=\vec{B'}}
\e{\vec{E}&=\vec{E'}+\fr{1}{c}\vec{B'}\times \vec{v}\\
&=\fr{1}{c}((B'_yv_z-B_z'v_y)\hat{x}-(B_x'v_z-B_z'v_x)\hat{y}+(B_x'v_y-B_y'v_x)\hat{z})
\intertext{$B_x'=B_x=0$ and $B_y'=B_y=0$.}
&=\fr{1}{c}(-B_z'v_y\hat{x}+B_z'v_x\hat{y}+0\hat{z})
\intertext{$E_x=0$.}
&=\label{iundia}\fr{1}{c}(0\hat{x}+B_z'v_x\hat{y}+0\hat{z})}
\eref{iundia} determines $v_y$.
\e{0=-\fr{1}{c}B_z'v_y}
Because $B_z'\neq 0$, $v_y=0$. Because $E_y\neq 0$,
\e{v_x=\fr{cE_y}{B_z'}}
The magnetic field in the boosted frame is
\e{\boxed{\vec{B'}=0\hat{x}+0\hat{y}+B_z'\hat{z}=0\hat{x}+0\hat{y}+\fr{c E_y}{v_x}\hat{z}}}
To find $v_z$, observe that the electrons are in the xy plane, so $v_z=0$. The boost velocity vector is
\e{\boxed{\vec{v}=\fr{cE_y}{B_z}\hat{x}+0\hat{y}+0\hat{z}}}

\section{Question 2.b.}
The previous result solves for $\vec{v}$.
\e{\boxed{\vec{j}=ne\vec{v}=ne(\fr{cE_y}{B_z}\hat{x}+0\hat{y}+0\hat{z})}}
Because $\vec{E}$ is in the y direction, $\vec{j}$ is \fbox{perpendicular} to $\vec{E}$.

\section{Question 2.c.}
Previously,
\e{\vec{E}=0\hat{x}+\fr{1}{c}B_zv_x\hat{y}+0\hat{z}}
and
\e{\vec{j}=\fr{neE_yc}{B_z}\hat{x}+0\hat{y}+0\hat{z}}
so
\e{j_x=\sxx E_x + \sxy E_y = \sxy E_y = \fr{nec}{B_z}E_y}
\e{\sxy=\fr{nec}{B_z}}
\e{\syx=\fr{-nec}{B_z}}
\e{j_y=\syx E_x + \syy E_y = \syy E_y = 0}
\e{j_y=0\quad E_y\neq 0 \Rightarrow \syy=0}
\e{\boxed{\sxx=0\quad \syy=0}}
\e{\boxed{\sxy=\fr{nec}{B_z}\quad \syx=-\fr{nec}{B_z}}}

\section{Question 2.d.}
Redo the previous parts using a Lorentz transformation. The only assumptions are $\vec{E}$, $\vec{B}$, and $\vec{E'}=0$.
\e{\vec{E}=0\hat{x}+E_y\hat{y}+0\hat{z}}
\e{\vec{B}=0\hat{x}+0\hat{y}+B_z\hat{z}}
The undetermined components are $E_y$ and $E_z$.
\e{E_y'=\lorentz (E_y+\fr{v}{c}B_z)}
\e{E_y=-\fr{v}{c}B_z}
\e{E_z'=\lorentz (E_z+\fr{v}{c}B_y)=0}
\e{B_x'=B_x=0}
\e{B_y'=\lorentz (B_y-\fr{v}{c}E_z)=0}
\e{B_y=\fr{v}{c}E_z}
\e{B_z'&= \lorentz (B_z+\fr{v}{c}E_y)\\
&=\lorentz (B_z-\fr{v^2}{c^2}B_z)\\
&=\lorentz B_z}
\e{\vec{j}=ne\vec{v}==ne\fr{cE_y}{B_z}\hat{x}+0\hat{y}+0\hat{z}}
\e{j_x&=\sxx E_x + \sxy E_y\\
&=\sxy E_y\\
&= -\fr{nec}{B_z}E_y}
\e{\sxy = - \fr{nec}{B_z}}
\e{j_y=\syx E_x + \syy E_y = \syy E_y = 0}
\e{E_y=0\Rightarrow \syy=0}
\e{\sxx=0}
The resistivity tensor is
\e{\boxed{\sxx = 0 \quad \syy = 0 \quad \sxy = -\fr{nec}{B_z} \quad \syx = \fr{nec}{Bz}}}
The $\vec{B'}$ field is
\e{\boxed{\vec{B'}=0\hat{x}+0\hat{y}+\lorentz B_z \hat{z}}}
The resistivity tensor $\fbox{does change}$.

\section{Question 3.a.}
\n{deltdtphi}{\delta \dtphi}
\n{deltdxphi}{\delta d_x \phi}
\n{deltphi}{\delta \phi}
\n{dtphi}{d_t \phi}
\n{dxphi}{d_x \phi}
\n{lag}{\mathcal{L}}
\n{oocs}{\fr{1}{c^2}}
\n{padtphi}{\fr{\partial}{\partial \dtphi}}
\n{padt}{\fr{\partial}{\partial t}}
\n{padxphi}{\fr{\partial}{\partial d_x \phi}}
\n{paphi}{\fr{\partial}{\partial \phi}}
\n{patsq}{\fr{\partial^2}{\partial t^2}}
\n{paxsq}{\fr{\partial^2}{\partial x^2}}
\n{pax}{\fr{\partial}{\partial x}}
\n{pat}{\fr{\partial}{\partial t}}
\n{phio}{\phi_1}
\n{phit}{\phi_2}
\n{dt}{\fr{d}{dt}}
\n{dx}{\fr{d}{dx}}
\n{dtsq}{d_t^2}
\n{dxsq}{d_x^2}
\n{notsurface}{\int (\deltphi \paphi \lag - \deltphi \dt \padtphi \lag - \deltphi \dx \padxphi \lag) dx dt}
\n{surface}{\int \padtphi (\padtphi \lag)_{t_1}^{t_2}dx + \int \deltdxphi (\padxphi \lag)_{t_1}^{t_2} dx}
\n{eom}{\paphi \lag - \dt \padtphi \lag - \dx \padxphi \lag}
On notation, $d_x$ is the total derivative, $\deltphi$ is the perturbation of $\phi$, $\deltdxphi$ is the perturbation of $\dx \phi$, and so on.
\e{S_1=\int \lag(\phi, \dtphi, \dxphi) dx dt}
Perturb.
\e{S_2&=\int \lag(\phi + \deltphi, \dtphi +\deltdtphi, \dxphi + \deltdxphi) dx dt\\
&=S_1+\int (\deltphi \paphi \lag + \deltdtphi \padtphi \lag + \deltdxphi \padxphi \lag) dx dt\\
\intertext{Integrate by parts.
\e{\int \deltdtphi \padtphi \lag dx dt=\int \deltphi (\padtphi \lag)_{t_1}^{t_2} dx - \int \deltphi \dt \pat \lag dx dt}
\e{\int \deltdxphi \padxphi \lag dx dt=\int \deltphi (\padxphi \lag)_{t_1}^{t_2} dx - \int \deltphi \dx \pax \lag dx dt}}
&=S_1 + \notsurface + \surface
\intertext{Set the ``surface terms'' to 0.
\e{\surface=0}}
&=S_1 + \notsurface\\
\intertext{Factor out $\deltphi$.}
&=S_1 + \int \deltphi (\eom) dx dt}
The perturbed action gives the unperturbed action, plus a differential equation in $\lag$. It is the Equation of Motion.
\e{\boxed{\eom=0}}
Next, substitute $\lag$ into the EOM to find the EOM in terms of $\phi$.
\e{\lag = \fr{1}{2c^2}(\dtphi)^2-\fr{1}{2}(\dxphi)^2-\fr{1}{2}m^2\phi^2}
\e{\boxed{m^2\phi-\fr{1}{c^2}\dtsq\phi+\dxsq\phi=0}}
Next, test if the linear combination $\phio+\phit$ solves the EOM.
\e{\phi=\phio+\phit}
\e{-m^2\phio-m^2\phit-\fr{1}{c^2}\dtsq(\phio+\phit)+\dxsq(\phio+\phit)=0}
\e{(-m^2\phio-\fr{1}{c^2}\dtsq\phio+\dxsq\phio)+(-m^2\phit-\fr{1}{c^2}\dtsq\phit+\dxsq\phit)=0}
The linear combination splits into two EOMs, each of which are solved.
\e{\boxed{0+0=0}}
In regards to the minutia, the integration by parts enacts a total derivative in $x$ and $t$, not a partial derivative, which gives the $\dxsq \phi$ and $\dtsq \phi$ terms in the EOM, instead of something like $\pax \dx \phi$ which is incorrect.

\section{Question 3.b.}
Substitute $\lag$ into the EOM to find the EOM in terms of $\phi$.
\e{\lag={\fr{1}{2c^2}\dtphi^2-\fr{1}{2}(\dxphi)^2 - \fr{1}{2} m^2 \phi^2 - \fr{\lambda}{4} \phi^4}}
\e{\label{91283n}\boxed{-m^2\phi-\fr{1}{c^2}\dtsq \phi+\dxsq \phi - \lambda \phi^3=0}}
Test if the sum of two solutions is also a solution.
\e{\phi=\phio+\phit}
\e{0&=-m^2\phio-m^2\phit-\oocs \dtsq \phio - \oocs \dtsq \phit + \dxsq \phio + \dxsq \phit - \lambda(\phio + \phit)^3\\
&=(-m^2\phio-\oocs \dtsq \phio+ \dxsq \phio)+(-m^2\phit- \oocs \dtsq \phit+ \dxsq \phit)-\lambda(\phio^3+\phit^3+3\phio^2\phit+3\phit^2\phio)\\
&=-3\lambda(\phio^2\phit+\phit^2\phio)}
\e{\label{ap9w8rhq}\boxed{-3\lambda(\phio^2\phit+\phit^2\phio)=0}}
In order for the equations of motion \eref{91283n} to be satisfied, \eref{ap9w8rhq} must also be satisfied, which means the field $\phi=\phio+\phit$ is \fbox{not necessarily a solution} to the EOM. As a comment, these could be called ``secondary'' equations of motion, and it would be interesting to find out if they are necessary to correctly predict dynamics, especially for theories such as $\lambda \phi^4$.

\end{document}